\documentclass[12pt,a4paper]{article}
\usepackage[utf8]{inputenc}
\usepackage[portuguese]{babel}
\usepackage{amsmath}
\usepackage{amssymb}
\usepackage{booktabs}
\usepackage{listings}
\usepackage{xcolor}
\usepackage{geometry}
\usepackage{hyperref}

\geometry{top=3cm, bottom=3cm, left=2.5cm, right=2.5cm}

% Configuração para exibição de código C++
\definecolor{codegreen}{rgb}{0,0.6,0}
\definecolor{codegray}{rgb}{0.5,0.5,0.5}
\definecolor{codepurple}{rgb}{0.58,0,0.82}
\definecolor{backcolour}{rgb}{0.95,0.95,0.92}

\lstdefinestyle{mystyle}{
    backgroundcolor=\color{backcolour},   
    commentstyle=\color{codegreen},
    keywordstyle=\color{magenta},
    numberstyle=\tiny\color{codegray},
    stringstyle=\color{codepurple},
    basicstyle=\ttfamily\footnotesize,
    breakatwhitespace=false,         
    breaklines=true,                 
    captionpos=b,                    
    keepspaces=true,                 
    numbers=left,                    
    numbersep=5pt,                  
    showspaces=false,                
    showstringspaces=false,
    showtabs=false,                  
    tabsize=2
}
\lstset{style=mystyle}

\title{Relatório Técnico: Resolução do Problema da Mochila 0-1 usando CPLEX Concert C++}
\author{Caio Anderson F.}
\date{\today}

\begin{document}

\maketitle

\section{Introdução}
O Problema da Mochila 0-1 (0-1 Knapsack Problem) é um clássico problema de otimização combinatória NP-difícil. Dado um conjunto de itens, cada um com um peso e um valor (lucro), o objetivo é selecionar um subconjunto de itens de forma a maximizar o valor total acumulado sem que a soma dos pesos dos itens selecionados ultrapasse a capacidade máxima da mochila.

\section{Formulação Matemática}
Seja $N$ o número total de itens disponíveis e $W$ a capacidade máxima da mochila. Para cada item $i \in \{1, \dots, N\}$, define-se:
\begin{itemize}
    \item $v_i$: o valor (lucro) do item $i$ ($v_i \ge 0$).
    \item $w_i$: o peso do item $i$ ($w_i \ge 0$).
\end{itemize}

O problema pode ser formulado como um modelo de Programação Linear Inteira Binária (PLIB) através da definição das variáveis de decisão:
\[
x_i = \begin{cases} 
1, & \text{se o item } i \text{ for selecionado para a mochila} \\
0, & \text{caso contrário}
\end{cases}
\]

A formulação matemática é expressa por:

\begin{alignat}{2}
\text{Maximizar} \quad & \sum_{i=1}^{N} v_i x_i \label{eq:obj} \\
\text{Sujeito a} \quad & \sum_{i=1}^{N} w_i x_i \le W \label{eq:capacity} \\
& x_i \in \{0, 1\}, \quad \forall i \in \{1, \dots, N\} \label{eq:binary}
\end{alignat}

Onde:
\begin{itemize}
    \item A Equação \eqref{eq:obj} define a função objetivo, maximizando o lucro acumulado dos itens selecionados.
    \item A Restrição \eqref{eq:capacity} garante que a soma dos pesos dos itens inseridos na mochila não exceda a capacidade limite $W$.
    \item A Restrição \eqref{eq:binary} impõe que a variável de escolha seja binária (decisão 0-1).
\end{itemize}

\newpage
\section{Implementação Computacional em C++}
A seguir é apresentado o trecho principal do solver implementado no arquivo \texttt{mochila\_solver.cpp}, utilizando a API \texttt{CPLEX Concert} em C++:

\begin{lstlisting}[language=C++, caption={Classe KnapsackSolver com formulação CPLEX Concert}]
class KnapsackSolver {
public:
    static void solve(const KnapsackInstance& instance) {
        int N = instance.get_N();
        double W = instance.get_capacity();
        IloEnv env;

        try {
            IloModel model(env);

            // Variaveis de decisao: x[i] = 1 se selecionado, 0 caso contrario
            IloBoolVarArray x(env, N);

            // Funcao Objetivo: Maximizar o lucro total
            IloExpr obj_expr(env);
            for (int i = 0; i < N; ++i) {
                obj_expr += instance.get_value(i) * x[i];
            }
            model.add(IloMaximize(env, obj_expr));
            obj_expr.end();

            // Restricao de Capacidade da Mochila: peso total <= W
            IloExpr weight_expr(env);
            for (int i = 0; i < N; ++i) {
                weight_expr += instance.get_weight(i) * x[i];
            }
            model.add(weight_expr <= W);
            weight_expr.end();

            IloCplex cplex(model);
            cplex.setOut(env.getNullStream()); // Oculta o log padrao do CPLEX
            cplex.setParam(IloCplex::Param::Threads, 1);

            // Exportacao do modelo matematico para depuracao em formato LP
            cplex.exportModel("knapsack_model.lp");

            auto start_time = std::chrono::high_resolution_clock::now();
            if (!cplex.solve()) {
                std::cerr << "Modelo inviavel ou falha no solver." << std::endl;
                env.end();
                return;
            }
            auto end_time = std::chrono::high_resolution_clock::now();
            double elapsed = std::chrono::duration<double>(end_time - start_time).count();

            // Recuperacao e exibicao de estatisticas finais
            double total_value = cplex.getObjValue();
            double total_weight = 0.0;
            std::vector<int> selected_items;
            for (int i = 0; i < N; ++i) {
                if (cplex.getValue(x[i]) > 0.5) {
                    selected_items.push_back(i);
                    total_weight += instance.get_weight(i);
                }
            }
            // Exibicao dos resultados omitida para concisao...
        } catch (IloException& e) {
            std::cerr << "Erro no Concert: " << e << std::endl;
        }
        env.end();
    }
};
\end{lstlisting}

\newpage
\section{Formatos de Instâncias Suportados}
O resolvedor foi dotado de um parser robusto que realiza a detecção explícita e automática de dois tipos comuns de instâncias:

\subsection{Formato de Cabeçalho Estruturado (Chave-Valor)}
Este formato é caracterizado por metadados descritos explicitamente e um marcador obrigatório delimitando o início da listagem de itens (\texttt{ITEMS} ou \texttt{DATA}). Permite ainda nomes de itens contendo espaços:
\begin{verbatim}
N: 4
CAPACITY: 100
ITEMS
Item Super Premium A 50 40
Item_B 30 20
\end{verbatim}

\subsection{Formato Padrão da Literatura (KPLIB)}
Compatível com as instâncias oficiais da \texttt{KPLIB} (Knapsack Problem Library), sem strings de marcação ou delimitadores chave-valor. O arquivo é composto unicamente por inteiros e reais:
\begin{verbatim}
50
14778
845 804
758 448
\end{verbatim}

\section{Resultados Computacionais}
O solver foi submetido a baterias de testes com instâncias reais contendo 1.000 itens retiradas da biblioteca benchmark \texttt{KPLIB}. Os tempos de execução refletem o poder de otimização combinatória do resolvedor com CPLEX em hardware moderno:

\begin{table}[htbp]
\centering
\caption{Resultados dos testes computacionais na base KPLIB ($N=1000$)}
\label{tab:results}
\begin{tabular}{lcccc}
\toprule
\textbf{Categoria da Instância} & \textbf{Lucro Ótimo} & \textbf{Capacidade Utilizada / Limite} & \textbf{Tempo (s)} \\
\midrule
Uncorrelated & 400.792 & 252.461 / 252.480 & 0,0096 s \\
Weakly Correlated & 273.039 & 245.968 / 245.972 & 0,0140 s \\
Strongly Correlated & 316.372 & 245.972 / 245.972 & 0,0035 s \\
Inverse Strongly Correlated & 263.977 & 295.477 / 295.477 & 0,0032 s \\
Almost Strongly Correlated & 316.408 & 245.972 / 245.972 & 0,0049 s \\
Subset Sum & 245.972 & 245.972 / 245.972 & 0,0025 s \\
Uncorrelated Similar Weights & 371.246 & 49.525.300 / 49.530.200 & 0,0025 s \\
Spanner Uncorrelated & 228.624 & 84.656 / 84.660 & 0,0009 s \\
Spanner Weakly Correlated & 196.050 & 198.664 / 198.695 & 0,0015 s \\
Spanner Strongly Correlated & 488.704 & 198.804 / 198.807 & 0,0009 s \\
Multiple Strongly Correlated & 399.170 & 245.970 / 245.970 & 0,0045 s \\
Profit Ceiling & 245.874 & 245.971 / 245.972 & 0,0025 s \\
Circle & 5.182.850 & 245.972 / 245.972 & 0,0025 s \\
\bottomrule
\end{tabular}
\end{table}

\section{Repositórios e Recursos}
Os códigos-fonte desenvolvidos e as bases de instâncias utilizadas neste trabalho estão disponíveis nos seguintes endereços eletrônicos:
\begin{itemize}
    \item \textbf{Código-Fonte em C++ (CPLEX)}: Repositório com a implementação completa do resolvedor, instruções de compilação e exemplos de execução. \\
    \url{https://github.com/caioandersonf/cplex-aplicado-ao-problema-da-mochila}
    \item \textbf{Biblioteca de Instâncias (KPLIB)}: Repositório contendo o conjunto completo de instâncias de benchmark da literatura do problema da mochila 0-1. \\
    \url{https://github.com/likr/kplib}
\end{itemize}

\end{document}
